\documentclass[11pt,a4paper]{article}
\usepackage[utf8]{inputenc}\usepackage[T1]{fontenc}\usepackage{lmodern}
\usepackage{geometry}\usepackage{amsmath,amssymb}\usepackage{xcolor}
\usepackage{listings}\usepackage{hyperref}\usepackage{microtype}
\geometry{margin=1in}\hypersetup{hidelinks}
\lstset{basicstyle=\footnotesize\ttfamily,breaklines=true,columns=fullflexible,
  frame=leftline,framesep=6pt,xleftmargin=8pt,aboveskip=5pt,belowskip=5pt}
\definecolor{axc}{HTML}{B8860B}\definecolor{onaxc}{HTML}{3A6EA5}
\definecolor{provedc}{HTML}{3C8A3C}
\newcommand{\tagax}{\textcolor{axc}{\textbf{[imported]}}}
\newcommand{\tagon}{\textcolor{onaxc}{\textbf{[rests on an import]}}}
\newcommand{\tagpr}{\textcolor{provedc}{\textbf{[proved outright]}}}
\title{Blueprint: what the Lean development proves}
\author{Generated by \texttt{scripts/blueprint.py}}\date{\today}
\begin{document}\maketitle

\noindent
Generated from the formalization, not written alongside it. Dependencies
are read from Lean's own printed proof terms; the classification is
propagated from the imported axioms and agrees with \verb|#print axioms|.
Signatures are transliterated to ASCII, and the generator aborts on any
character it has no mapping for rather than substituting a placeholder.

\medskip\noindent
\tagax\ assumed. Both are classical theorems, neither open nor near the
frontier this development concerns.\\
\tagon\ proved, but its proof reaches an imported axiom.\\
\tagpr\ proved from nothing beyond Lean's own three axioms.

\medskip\noindent
Of 57 declarations, 45 are proved outright, 10 rest on an
import, and 2 are imported. No \verb|sorry| occurs anywhere.


\section{Imported results}
\noindent\texttt{Diaz/Axioms.lean}


\subsection*{\texttt{exists\_ringHom\_of\_transcendental} \hfill \tagax}

**Steinitz extension.** If `u` and `t` are both transcendental over a subfield `K` of `C`, some ring endomorphism of `C` fixes `K` pointwise and sends `u` to `t`. Why it is true: `u |-> t` is an isomorphism `K(u) -> K(t)` fixing `K`, because both are transcendental, so both are rational function fields over `K`. Extend a transcendence basis of `C` over `K(u)` into one over `K(t)` --- both have the cardinality of the continuum --- and then extend algebraically, `C` being algebraically closed. See Lang, *Algebra*, 3rd ed., Ch. VIII, or Steinitz's theorem in any account of field theory. The thesis is deliberately weak: a ring *endomorphism*, not an automorphism, which is all the results below consume, and which is automatically injective since `C` is a field. The hypothesis is necessary --- for `u` algebraic over `K` and `t` not, no such map exists. Mathlib has the ingredients (`IsAlgClosed.equivOfTranscendenceBasis`, `IsAlgClosed.lift`) but not the assembled statement.

\begin{lstlisting}
axiom exists_ringHom_of_transcendental {K : Subfield C} {u t : C}
 (hu : Transcendental (K) u) (ht : Transcendental (K) t) :
 exists Phi : C ->+* C, (forall a in K, Phi a = a) and Phi u = t

end Diaz
\end{lstlisting}

\subsection*{\texttt{hermite\_lindemann} \hfill \tagax}

**Hermite--Lindemann.** If `u != 0` and `exp u` is algebraic, then `u` is transcendental. This is the contrapositive of the usual statement, that a non-zero algebraic number has transcendental exponential. See A. Baker, *Transcendental Number Theory*, Cambridge University Press 1975, Theorem 1.4. **This will become dischargeable.** Mathlib master carries only the analytic half (`NumberTheory/Transcendental/Lindemann/AnalyticalPart`), but the theorem itself is formalized in the open PR `leanprover-community/mathlib4\#28013`, as theorem transcendental\_exp \{a : C\} (a0 : a != 0) (ha : IsAlgebraic Z a) : Transcendental Z (exp a) which is the contrapositive of this axiom, over `Z` rather than `Q` -- the same condition in characteristic zero. When that PR merges, delete this axiom and derive it; the statement below is shaped to make that a local change.

\begin{lstlisting}
axiom hermite_lindemann {u : C} (hu : u != 0)
 (hexp : IsAlgebraic Q (Complex.exp u)) : Transcendental Q u

/- 
 
 

 
 
 
 
 
 

 
 
 
 

 
 -/
\end{lstlisting}

\section{The closure theorem}
\noindent\texttt{Diaz/Closure.lean}


\subsection*{\texttt{conj\_eq\_rho\_div} \hfill \tagpr}

With `rho = u * conj u`, the conjugate of `u` is `rho / u`. Trivial as algebra, and it is the entire content of the closure theorem: `conj u` is not an independent quantity but a rational function of `u` over the base field. Everything else follows.

\begin{lstlisting}
theorem conj_eq_rho_div (hu : u != 0) : (u * conj u) / u = conj u
\end{lstlisting}

\subsection*{\texttt{conj\_mem\_hull} \hfill \tagpr}

Complex conjugation maps `hull K u` into itself, given that `K` is conjugation-stable and contains `rho = u * conj u`.

\begin{lstlisting}
theorem conj_mem_hull (hK : forall z in K, conj z in K) (hu : u != 0)
 (hrho : u * conj u in K) :
 forall z in hull K u, conj z in hull K u
\end{lstlisting}
\noindent\footnotesize Uses: \texttt{conj\_eq\_rho\_div}, \texttt{mem\_hull\_of\_mem\_base}, \texttt{self\_mem\_hull}\normalsize


\subsection*{\texttt{conj\_not\_linear} \hfill \tagpr}

Conjugation is **not** `K`-linear once `K` contains a non-real number. A `K`-linear map `f` would satisfy `f (a * z) = a * f z` for `a in K`; taking `z = 1` forces `conj a = a`, so only a totally real base field could carry conjugation as a linear map. Since the intended base is the algebraic numbers, which contain `i`, conjugation on the hull is a ring involution and nothing stronger.

\begin{lstlisting}
theorem conj_not_linear {a : C} (ha : conj a != a) :
 not (forall z : C, conj (a * z) = a * conj z)
\end{lstlisting}

\subsection*{\texttt{conj\_not\_linear\_hull} \hfill \tagpr}

The form actually needed: conjugation is not `K`-linear *on the hull*. The refuted claim is about the hull, and `forall z in hull` is weaker than `forall z : C`, so its negation is the stronger statement. It is free: the witness `z = 1` lies in the hull.

\begin{lstlisting}
theorem conj_not_linear_hull {a : C} (hne : conj a != a) :
 not (forall z in hull K u, conj (a * z) = a * conj z)
\end{lstlisting}
\noindent\footnotesize Uses: \texttt{hull}\normalsize


\subsection*{\texttt{conj\_not\_linear\_of\_I} \hfill \tagpr}

Concretely: conjugation is not linear over any base containing `i`.

\begin{lstlisting}
theorem conj_not_linear_of_I :
 not (forall z : C, conj (Complex.I * z) = Complex.I * conj z)
\end{lstlisting}
\noindent\footnotesize Uses: \texttt{conj\_not\_linear}\normalsize


\subsection*{\texttt{eqOn\_hull} \hfill \tagpr}

Two ring homomorphisms of `C` agreeing on `K` and at `u` agree on the whole hull. This is the "a homomorphism is determined by these data" step of the closure theorem's proof.

\begin{lstlisting}
theorem eqOn_hull (f g : C ->+* C) (hK : forall z in K, f z = g z) (hu : f u = g u) :
 forall z in hull K u, f z = g z
\end{lstlisting}
\noindent\footnotesize Uses: \texttt{hull}\normalsize


\subsection*{\texttt{hull} \hfill \tagpr}

The subfield of `C` generated by `K` together with `u`.

\begin{lstlisting}
noncomputable def hull : Subfield C
\end{lstlisting}

\subsection*{\texttt{mem\_hull\_of\_mem\_base} \hfill \tagpr}
\begin{lstlisting}
theorem mem_hull_of_mem_base {z : C} (hz : z in K) : z in hull K u
\end{lstlisting}
\noindent\footnotesize Uses: \texttt{hull}\normalsize


\subsection*{\texttt{ne\_zero\_of\_notMem} \hfill \tagpr}

An element outside a subfield is non-zero, subfields containing `0`.

\begin{lstlisting}
theorem ne_zero_of_notMem {K : Subfield C} {z : C} (h : z notin K) : z != 0
\end{lstlisting}

\subsection*{\texttt{self\_mem\_hull} \hfill \tagpr}
\begin{lstlisting}
theorem self_mem_hull : u in hull K u
\end{lstlisting}
\noindent\footnotesize Uses: \texttt{hull}\normalsize


\subsection*{\texttt{transcendental\_candidate\_over\_base} \hfill \tagon}

The bridge, assembled: a candidate is transcendental over the base.

\begin{lstlisting}
theorem transcendental_candidate_over_base {L : Subfield C}
 [Algebra.IsAlgebraic Q (L)] (hu : u != 0)
 (hexp : IsAlgebraic Q (Complex.exp u)) : Transcendental (L) u
\end{lstlisting}
\noindent\footnotesize Uses: \texttt{transcendental\_of\_base}, \texttt{transcendental\_of\_candidate}\normalsize


\subsection*{\texttt{transcendental\_ne\_zero} \hfill \tagpr}

A transcendental element is non-zero: `0` is a root of `X`.

\begin{lstlisting}
theorem transcendental_ne_zero {F : Type*} [Field F] [Algebra F C] {z : C}
 (h : Transcendental F z) : z != 0
\end{lstlisting}

\subsection*{\texttt{transcendental\_of\_base} \hfill \tagpr}

Transcendence over `Q` upgrades to transcendence over any base algebraic over `Q` --- in particular over the algebraic numbers, which is the intended base. Without this the imported axiom is a dead leaf: the results above take transcendence over the base as a hypothesis, while Hermite--Lindemann supplies it only over `Q`. The algebraicity hypothesis is necessary rather than decorative: for a base containing `u` the conclusion is false.

\begin{lstlisting}
theorem transcendental_of_base {L : Subfield C} [Algebra.IsAlgebraic Q (L)]
 {z : C} (h : Transcendental Q z) : Transcendental (L) z
\end{lstlisting}

\subsection*{\texttt{transcendental\_of\_candidate} \hfill \tagon}

A counterexample to Diaz's conjecture is transcendental. The only imported input is Hermite--Lindemann.

\begin{lstlisting}
theorem transcendental_of_candidate (hu : u != 0)
 (hexp : IsAlgebraic Q (Complex.exp u)) : Transcendental Q u
\end{lstlisting}
\noindent\footnotesize Uses: \texttt{hermite\_lindemann}\normalsize


\section{The model}
\noindent\texttt{Diaz/Model.lean}


\subsection*{\texttt{conj\_coords} \hfill \tagpr}

Complex conjugation acts on the plane by swapping coordinates. This is what licenses stating the involution in coordinates.

\begin{lstlisting}
theorem conj_coords (t : C) (a b : Q) :
 conj ((a : C) * t + (b : C) * conj t) = (b : C) * t + (a : C) * conj t
\end{lstlisting}

\subsection*{\texttt{exists\_transcendental\_on\_circle} \hfill \tagon}

**A transcendental point on the circle.** For any non-zero `r` in a base `K` algebraic over `Q`, the number `t = r * e\^{}i` is transcendental over `K`, with `t * conj t` equal to `r * conj r` **and** in the base. Both conclusions are needed: the model lemmas consume the membership, while the transfer results consume the equality `t * conj t = u * conj u`. An earlier version gave only the membership, so no exhibited `t` could discharge the transfer hypotheses. The hypothesis `conj r in L` is needed and is not implied by `r in L`: for `L = Q(2\^{}\{1/3\}omega)` and `r = 2\^{}\{1/3\}omega` one has `r * conj r = 2\^{}\{2/3\} notin L`. Every result that consumes this already assumes `L` conjugation-stable. So the hypotheses of `norm\_mem\_iff` are not merely satisfiable in the abstract; they are met by an explicit complex number.

\begin{lstlisting}
theorem exists_transcendental_on_circle {L : Subfield C}
 [Algebra.IsAlgebraic Q (L)] {r : C} (hr : r in L) (hrc : conj r in L)
 (hr0 : r != 0) :
 exists t : C, t != 0 and Transcendental (L) t and t * conj t = r * conj r
 and t * conj t in L
\end{lstlisting}
\noindent\footnotesize Uses: \texttt{transcendental\_exp\_I}, \texttt{transcendental\_of\_base}\normalsize


\subsection*{\texttt{exists\_transcendental\_on\_circle\_sq} \hfill \tagon}

The circle condition in the form the matrix results consume. `exists\_transcendental\_on\_circle` gives `t * conj t = r * conj r`, while `Rigidity` states its hypotheses as `u * conj u = r \^{} 2`. These agree exactly when `r` is real --- the intended case, `r = sqrt rho` with `rho = |u|\^{}2` --- which is not implied by `r in L` and so has to be said.

\begin{lstlisting}
theorem exists_transcendental_on_circle_sq {L : Subfield C}
 [Algebra.IsAlgebraic Q (L)] {r : C} (hr : r in L) (hrr : conj r = r)
 (hr0 : r != 0) :
 exists t : C, t != 0 and Transcendental (L) t and t * conj t = r ^ 2
 and t * conj t in L
\end{lstlisting}
\noindent\footnotesize Uses: \texttt{exists\_transcendental\_on\_circle}\normalsize


\subsection*{\texttt{indep} \hfill \tagpr}

`t` and `conj t` are independent over `Q`. This is what makes coordinates `(a, b) |-> a t + b conj t` well defined, so that a function of the coordinates is a function on the plane. It is the analogue of the axis lemma.

\begin{lstlisting}
theorem indep (hT : Transcendental K t) (hrho : t * conj t in K)
 {a b : Q} (h : (a : C) * t + (b : C) * conj t = 0) : a = 0 and b = 0
\end{lstlisting}
\noindent\footnotesize Uses: \texttt{conj\_eq\_rho\_div}, \texttt{sq\_add\_sq\_notMem}, \texttt{transcendental\_ne\_zero}\normalsize


\subsection*{\texttt{norm\_form} \hfill \tagpr}

On the rational plane spanned by `t` and `conj t`, (a t + b tbar)(a tbar + b t) = (a\^{}2 + b\^{}2)*t tbar + a b*(t\^{}2 + tbar\^{}2). This is `x sigma(x) = (a\^{}2 + b\^{}2) rho + a b (T\^{}2 + rho\^{}2/T\^{}2)` of `prop:model`, with `conj` in place of `sigma`.

\begin{lstlisting}
theorem norm_form (t : C) (a b : Q) :
 ((a : C) * t + (b : C) * conj t) * conj ((a : C) * t + (b : C) * conj t)
 = ((a : C) ^ 2 + (b : C) ^ 2) * (t * conj t)
 + ((a : C) * (b : C)) * (t ^ 2 + (conj t) ^ 2)
\end{lstlisting}

\subsection*{\texttt{norm\_mem\_iff} \hfill \tagpr}

**The characterisation of `D0`.** For `a, b` rational, the element `x = a t + b conj t` has `x * conj x in K` exactly when `a = 0` or `b = 0`. So the elements of the rational plane with "algebraic modulus" are precisely the rational multiples of `t` and of `conj t` --- the two rays of `prop:model`, and in particular there are non-zero ones. Every algebraic constraint a Diaz candidate satisfies is satisfied here too.

\begin{lstlisting}
theorem norm_mem_iff (hT : Transcendental K t) (hrho : t * conj t in K) (a b : Q) :
 ((a : C) * t + (b : C) * conj t) * conj ((a : C) * t + (b : C) * conj t) in K
 <-> a = 0 or b = 0
\end{lstlisting}
\noindent\footnotesize Uses: \texttt{norm\_form}, \texttt{sq\_add\_sq\_notMem}\normalsize


\subsection*{\texttt{sq\_add\_sq\_notMem} \hfill \tagpr}

If `t` is transcendental over `K` and `rho = t * conj t` lies in `K`, then the cross term `t\^{}2 + conj(t)\^{}2` does not lie in `K`. Otherwise `t` would satisfy `X\^{}4 - c X\^{}2 + rho\^{}2`, a non-zero polynomial over `K`. This is what forces `a b = 0` in the model.

\begin{lstlisting}
theorem sq_add_sq_notMem (hT : Transcendental K t) (hrho : t * conj t in K) :
 t ^ 2 + (conj t) ^ 2 notin K
\end{lstlisting}

\subsection*{\texttt{transcendental\_exp\_I} \hfill \tagon}

`e\^{}i` is transcendental over `Q`. Immediate from Hermite--Lindemann, since `i` is algebraic and non-zero.

\begin{lstlisting}
theorem transcendental_exp_I : Transcendental Q (Complex.exp Complex.I)
\end{lstlisting}
\noindent\footnotesize Uses: \texttt{transcendental\_of\_candidate}\normalsize


\section{The formal exponential}
\noindent\texttt{Diaz/Exponential.lean}


\subsection*{\texttt{Exp0} \hfill \tagpr}

`Exp0 (a, b) = 2 \^{} (a + b)`, the formal exponential of the model.

\begin{lstlisting}
noncomputable def Exp0 (x : Q x Q) : R
\end{lstlisting}

\subsection*{\texttt{Exp0\_add} \hfill \tagpr}

`Exp0` is a homomorphism from the plane to `R\^{} x `.

\begin{lstlisting}
theorem Exp0_add (x y : Q x Q) : Exp0 (x + y) = Exp0 x * Exp0 y
\end{lstlisting}
\noindent\footnotesize Uses: \texttt{Exp0}\normalsize


\subsection*{\texttt{Exp0\_eq\_one\_iff} \hfill \tagpr}

**The kernel of this `Exp0` is the divisible line `a + b = 0`.** A property of this choice, not of the model --- see the caveat in the file header, where a variant with a lattice kernel is given.

\begin{lstlisting}
theorem Exp0_eq_one_iff (x : Q x Q) : Exp0 x = 1 <-> x.1 + x.2 = 0
\end{lstlisting}
\noindent\footnotesize Uses: \texttt{Exp0}, \texttt{two\_rpow\_eq\_one\_iff}\normalsize


\subsection*{\texttt{Exp0\_isAlgebraic} \hfill \tagpr}

Every value of `Exp0` is algebraic: the model really is a model of the *arithmetic* situation, not merely of the field structure.

\begin{lstlisting}
theorem Exp0_isAlgebraic (x : Q x Q) : IsAlgebraic Q (Exp0 x)
\end{lstlisting}
\noindent\footnotesize Uses: \texttt{Exp0}, \texttt{isAlgebraic\_two\_rpow}\normalsize


\subsection*{\texttt{Exp0\_ne\_zero} \hfill \tagpr}
\begin{lstlisting}
theorem Exp0_ne_zero (x : Q x Q) : Exp0 x != 0
\end{lstlisting}
\noindent\footnotesize Uses: \texttt{Exp0\_pos}\normalsize


\subsection*{\texttt{Exp0\_pos} \hfill \tagpr}
\begin{lstlisting}
theorem Exp0_pos (x : Q x Q) : 0 < Exp0 x
\end{lstlisting}
\noindent\footnotesize Uses: \texttt{Exp0}\normalsize


\subsection*{\texttt{Exp0\_pow\_eq\_one\_iff} \hfill \tagpr}

**The image of this `Exp0` is torsion-free.** A standard kernel `Zomega` makes `Exp (pomega/m)` a primitive `m`-th root of unity, so the image would have to contain every root of unity; `2\^{}Q` contains none but `1`. Again a property of this choice: the variant in the file header has both a lattice kernel and torsion in its image.

\begin{lstlisting}
theorem Exp0_pow_eq_one_iff (x : Q x Q) {n : N} (hn : n != 0) :
 Exp0 x ^ n = 1 <-> Exp0 x = 1
\end{lstlisting}
\noindent\footnotesize Uses: \texttt{Exp0\_eq\_one\_iff}\normalsize


\subsection*{\texttt{Exp0\_swap} \hfill \tagpr}

`Exp0` commutes with the involution. In coordinates the involution swaps `a` and `b`; the values are real, so complex conjugation fixes them, and `a + b` is symmetric.

\begin{lstlisting}
theorem Exp0_swap (x : Q x Q) : Exp0 (x.2, x.1) = Exp0 x
\end{lstlisting}
\noindent\footnotesize Uses: \texttt{Exp0}\normalsize


\subsection*{\texttt{Exp0\_swap\_conj} \hfill \tagpr}

The form the model actually needs: `Exp0 . sigma = conj . Exp0`, stated with the conjugation rather than only with the coordinate swap.

\begin{lstlisting}
theorem Exp0_swap_conj (x : Q x Q) :
 ((Exp0 (x.2, x.1) : C)) = (starRingEnd C) ((Exp0 x : C))
\end{lstlisting}
\noindent\footnotesize Uses: \texttt{Exp0\_swap}\normalsize


\subsection*{\texttt{isAlgebraic\_two\_rpow} \hfill \tagpr}

`2 \^{} q` is algebraic for every rational `q`: it is a root of `X \^{} q.den - 2 \^{} q.num`.

\begin{lstlisting}
theorem isAlgebraic_two_rpow (q : Q) : IsAlgebraic Q ((2 : R) ^ (q : R))
\end{lstlisting}

\subsection*{\texttt{model\_falsifies} \hfill \tagpr}

**The analogue of Diaz's conjecture is false in the model.** There is a non-zero element of the plane whose product with its conjugate lies in `K` and which carries an algebraic formal exponential. This is the claim the closure theorem is pointed at, and it is what makes the absence of a contradiction more than a failure to find one.

\begin{lstlisting}
theorem model_falsifies {K : Subfield C} {t : C}
 (hrho : t * conj t in K) (ht0 : t != 0) :
 exists a b : Q, ((a : C) * t + (b : C) * conj t) != 0
 and IsAlgebraic Q (Exp0 (a, b))
 and ((a : C) * t + (b : C) * conj t)
 * conj ((a : C) * t + (b : C) * conj t) in K
\end{lstlisting}
\noindent\footnotesize Uses: \texttt{Exp0\_isAlgebraic}\normalsize


\subsection*{\texttt{two\_rpow\_eq\_one\_iff} \hfill \tagpr}
\begin{lstlisting}
theorem two_rpow_eq_one_iff (r : R) : (2 : R) ^ r = 1 <-> r = 0
\end{lstlisting}

\section{Transfer}
\noindent\texttt{Diaz/Transfer.lean}


\subsection*{\texttt{Hmat\_transfer} \hfill \tagpr}

The transported matrix of a candidate is the matrix of its image: `Phi` carries `H u r` to `H t r` when it fixes `r` and sends `u` to `t`.

\begin{lstlisting}
theorem Hmat_transfer (Phi : C ->+* C) (hK : forall a in K, Phi a = a) {r : C} (hr : r in K)
 (hKconj : forall a in K, conj a in K) (hu0 : u != 0) (ht0 : t != 0) (hPhiu : Phi u = t)
 (hrho : u * conj u in K) (hrhot : t * conj t = u * conj u) :
 (Hmat u r).map Phi = Hmat t r
\end{lstlisting}
\noindent\footnotesize Uses: \texttt{Hmat}, \texttt{conj\_comm}, \texttt{self\_mem\_hull}\normalsize


\subsection*{\texttt{coeff\_transfer} \hfill \tagpr}

A ring homomorphism fixing `K` carries the coefficient `w\^{}T M v` to the coefficient of the transported matrix, for `w, v` over `K`. So vanishing of coefficients over `K` transfers.

\begin{lstlisting}
theorem coeff_transfer (Phi : C ->+* C) (hK : forall a in K, Phi a = a)
 {m n : N} (M : Matrix (Fin m) (Fin n) C) (w : Fin m -> C) (v : Fin n -> C)
 (hw : forall i, w i in K) (hv : forall j, v j in K) :
 Phi (sum i, sum j, w i * M i j * v j) = sum i, sum j, w i * (Phi (M i j)) * v j
\end{lstlisting}

\subsection*{\texttt{conj\_comm} \hfill \tagpr}

**An isomorphism matching the generators intertwines conjugation.** If `Phi` fixes `K` pointwise and sends `u` to `t`, where `u` and `t` both lie on the circle `z * conj z = rho` over `K`, then `Phi` commutes with complex conjugation on the whole hull of `u`. Note where this comes from: `conj u = rho / u` is a *rational function of `u` over `K`*, so conjugation is not extra data that could fail to be matched --- it is already determined by the ring structure. That is the content of the closure theorem.

\begin{lstlisting}
theorem conj_comm (Phi : C ->+* C) (hK : forall a in K, Phi a = a)
 (hKconj : forall a in K, conj a in K)
 (hu0 : u != 0) (ht0 : t != 0) (hPhiu : Phi u = t)
 (hrho : u * conj u in K) (hrhot : t * conj t = u * conj u) :
 forall z in hull K u, Phi (conj z) = conj (Phi z)
\end{lstlisting}
\noindent\footnotesize Uses: \texttt{conj\_eq\_rho\_div}, \texttt{eqOn\_hull}\normalsize


\subsection*{\texttt{exists\_conj\_intertwining} \hfill \tagon}

**The closure theorem, existence included.** Any two candidates over the same base, with the same `rho`, are carried onto one another by a ring endomorphism of `C` that fixes the base and intertwines conjugation. The intertwining is not arranged --- it follows, because `conj u = rho/u` makes conjugation a rational function of the generator.

\begin{lstlisting}
theorem exists_conj_intertwining (hKconj : forall a in K, conj a in K)
 (hu : Transcendental (K) u) (ht : Transcendental (K) t)
 (hrho : u * conj u in K) (hrhot : t * conj t = u * conj u) :
 exists Phi : C ->+* C, (forall a in K, Phi a = a) and Phi u = t
 and forall z in hull K u, Phi (conj z) = conj (Phi z)
\end{lstlisting}
\noindent\footnotesize Uses: \texttt{conj\_comm}, \texttt{exists\_ringHom\_of\_transcendental}, \texttt{transcendental\_ne\_zero}\normalsize


\subsection*{\texttt{hullsEquiv} \hfill \tagpr}

The hulls of any two elements transcendental over `K` are isomorphic as `K`-algebras, both being rational function fields.

\begin{lstlisting}
noncomputable def hullsEquiv (hu : Transcendental (K) u) (ht : Transcendental (K) t) :
 (IntermediateField.adjoin (K) {u}) ~=a [K] (IntermediateField.adjoin (K) {t})
\end{lstlisting}

\subsection*{\texttt{injective\_of\_hom} \hfill \tagpr}

A ring homomorphism of fields is injective. That injectivity is what *would* give preservation of `K`-linear independence, and hence of rank and structural rank; neither is formalized here, and this lemma states only the injectivity.

\begin{lstlisting}
theorem injective_of_hom (Phi : C ->+* C) : Function.Injective Phi
\end{lstlisting}

\section{The rank-one matrix}
\noindent\texttt{Diaz/Rigidity.lean}


\subsection*{\texttt{Hmat} \hfill \tagpr}

The matrix `H` attached to a candidate.

\begin{lstlisting}
noncomputable def Hmat (u r : C) : Matrix (Fin 2) (Fin 2) C
\end{lstlisting}

\subsection*{\texttt{candidate\_no\_vanishing\_coeff} \hfill \tagon}

If `u` is a Diaz candidate over a base algebraic over `Q`, then the attached rank-one matrix has no vanishing coefficient over that base.

\begin{lstlisting}
theorem candidate_no_vanishing_coeff {L : Subfield C} [Algebra.IsAlgebraic Q (L)]
 {u r : C} (hu0 : u != 0) (hexp : IsAlgebraic Q (Complex.exp u))
 (hr : r in L) (h : u * conj u = r ^ 2)
 (w v : Fin 2 -> C) (hwK : forall i, w i in L) (hvK : forall j, v j in L)
 (hw : w != 0) (hv : v != 0) :
 sum i, sum j, w i * (Hmat u r) i j * v j != 0
\end{lstlisting}
\noindent\footnotesize Uses: \texttt{no\_vanishing\_coeff\_matrix}, \texttt{transcendental\_candidate\_over\_base}\normalsize


\subsection*{\texttt{coeff\_eq\_matrix} \hfill \tagpr}
\begin{lstlisting}
theorem coeff_eq_matrix (u r : C) (w v : Fin 2 -> C) :
 sum i, sum j, w i * (Hmat u r) i j * v j
 = w 0 * (u * v 0 + r * v 1) + w 1 * (r * v 0 + conj u * v 1)
\end{lstlisting}
\noindent\footnotesize Uses: \texttt{Hmat}\normalsize


\subsection*{\texttt{coeff\_factor} \hfill \tagpr}

The coefficient `w\^{}T H v` factors, because `H` has rank one.

\begin{lstlisting}
theorem coeff_factor (hu0 : u != 0) (h : u * conj u = r ^ 2)
 (w v : Fin 2 -> C) :
 w 0 * (u * v 0 + r * v 1) + w 1 * (r * v 0 + conj u * v 1)
 = (w 0 + (r / u) * w 1) * (u * v 0 + r * v 1)
\end{lstlisting}

\subsection*{\texttt{det\_Hmat} \hfill \tagpr}

`det H = 0`: over `C` the rows are dependent.

\begin{lstlisting}
theorem det_Hmat (h : u * conj u = r ^ 2) : (Hmat u r).det = 0
\end{lstlisting}
\noindent\footnotesize Uses: \texttt{Hmat}\normalsize


\subsection*{\texttt{mem\_of\_lin\_rel} \hfill \tagpr}

If `u * a + r * b = 0` with `a, b` in `K`, not both zero, and `r != 0` in `K`, then `u in K`.

\begin{lstlisting}
theorem mem_of_lin_rel (hr : r in K) (hr0 : r != 0) {a b : C}
 (ha : a in K) (hb : b in K) (hab : not (a = 0 and b = 0))
 (h : u * a + r * b = 0) : u in K
\end{lstlisting}

\subsection*{\texttt{no\_vanishing\_coeff} \hfill \tagpr}

**No vanishing coefficient.** For non-zero `w, v` over `K`, the coefficient `w\^{}T H v` is non-zero.

\begin{lstlisting}
theorem no_vanishing_coeff (hr : r in K) (huK : u notin K)
 (h : u * conj u = r ^ 2)
 (w v : Fin 2 -> C) (hwK : forall i, w i in K) (hvK : forall i, v i in K)
 (hw : w != 0) (hv : v != 0) :
 w 0 * (u * v 0 + r * v 1) + w 1 * (r * v 0 + conj u * v 1) != 0
\end{lstlisting}
\noindent\footnotesize Uses: \texttt{coeff\_factor}, \texttt{mem\_of\_lin\_rel}, \texttt{ne\_zero\_of\_notMem}\normalsize


\subsection*{\texttt{no\_vanishing\_coeff\_matrix} \hfill \tagpr}

**No vanishing coefficient, in matrix form.**

\begin{lstlisting}
theorem no_vanishing_coeff_matrix (hr : r in K) (huK : u notin K)
 (h : u * conj u = r ^ 2)
 (w v : Fin 2 -> C) (hwK : forall i, w i in K) (hvK : forall j, v j in K)
 (hw : w != 0) (hv : v != 0) :
 sum i, sum j, w i * (Hmat u r) i j * v j != 0
\end{lstlisting}
\noindent\footnotesize Uses: \texttt{coeff\_eq\_matrix}, \texttt{no\_vanishing\_coeff}\normalsize


\section{The intended base}
\noindent\texttt{Diaz/Instantiation.lean}


\subsection*{\texttt{Qbar} \hfill \tagpr}

The algebraic numbers, as a subfield of `C`.

\begin{lstlisting}
noncomputable def Qbar : Subfield C
\end{lstlisting}

\subsection*{\texttt{QbarIsAlgebraic} \hfill \tagpr}

The instance that does not fire on its own. Named rather than anonymous so that it appears in the generated artefacts: an anonymous instance has no identifier for the generator to match, and would be invisible to the very document meant to expose the assumed surface.

\begin{lstlisting}
noncomputable instance QbarIsAlgebraic : Algebra.IsAlgebraic Q (Qbar)
\end{lstlisting}
\noindent\footnotesize Uses: \texttt{Qbar}\normalsize


\subsection*{\texttt{candidate\_indistinguishable} \hfill \tagon}

**A candidate is carried onto an ordinary complex number.** If `u` is a Diaz candidate over `Qbar` with `u * conj u = r\^{}2`, `r` real and algebraic, then there is a transcendental `t` and a ring endomorphism of `C` fixing `Qbar`, sending `u` to `t`, and intertwining conjugation on the hull of `u`. That `t` lies on the same circle is derivable from the statement, `Phi` fixing `Qbar` and the intertwining at `z = u`. The witness is `t = r * e\^{}i`, and it is *not* asserted that `exp t` is transcendental: were it algebraic, `t` would itself be a candidate, so that assertion is an instance of the conjecture rather than something proved here. This runs from the arithmetic hypotheses to the endomorphism with no step left in prose. It does not itself apply `coeff\_transfer` or `Hmat\_transfer`; those consume the endomorphism it produces.

\begin{lstlisting}
theorem candidate_indistinguishable
 {u r : C} (hu0 : u != 0) (hexp : IsAlgebraic Q (Complex.exp u))
 (hr : r in Qbar) (hrr : conj r = r) (hr0 : r != 0)
 (h : u * conj u = r ^ 2) :
 exists t : C, t != 0 and Transcendental (Qbar) t and 
 exists Phi : C ->+* C, (forall a in Qbar, Phi a = a) and Phi u = t
 and forall z in hull Qbar u, Phi (conj z) = conj (Phi z)
\end{lstlisting}
\noindent\footnotesize Uses: \texttt{conj\_mem\_Qbar}, \texttt{exists\_conj\_intertwining}, \texttt{exists\_transcendental\_on\_circle\_sq}, \texttt{transcendental\_candidate\_over\_base}\normalsize


\subsection*{\texttt{candidate\_no\_vanishing\_coeff\_Qbar} \hfill \tagon}

**The end-to-end statement over the intended base.** If `u` is a Diaz candidate --- non-zero, with `exp u` algebraic and `u * conj u = r\^{}2` for an algebraic `r` --- then the attached rank-one matrix has no vanishing coefficient over the algebraic numbers. This is `candidate\_no\_vanishing\_coeff` with `L = Qbar`, and its point is that the hypotheses are the arithmetic ones rather than transcendence, and that the base is the one the note is about rather than an arbitrary subfield.

\begin{lstlisting}
theorem candidate_no_vanishing_coeff_Qbar
 {u r : C} (hu0 : u != 0) (hexp : IsAlgebraic Q (Complex.exp u))
 (hr : r in Qbar) (h : u * conj u = r ^ 2)
 (w v : Fin 2 -> C) (hwK : forall i, w i in Qbar) (hvK : forall j, v j in Qbar)
 (hw : w != 0) (hv : v != 0) :
 sum i, sum j, w i * (Hmat u r) i j * v j != 0
\end{lstlisting}
\noindent\footnotesize Uses: \texttt{Qbar}, \texttt{candidate\_no\_vanishing\_coeff}\normalsize


\subsection*{\texttt{conj\_mem\_Qbar} \hfill \tagpr}

`Qbar` is conjugation-stable, which every result above assumes of the base. If `p` kills `a` then it kills `conj a`, its coefficients being rational and so fixed by conjugation.

\begin{lstlisting}
theorem conj_mem_Qbar {a : C} (h : a in Qbar) : conj a in Qbar
\end{lstlisting}
\noindent\footnotesize Uses: \texttt{mem\_Qbar\_iff}\normalsize


\subsection*{\texttt{exists\_transcendental\_on\_circle\_Qbar} \hfill \tagon}

And the circle carries a transcendental point over that base, so the model results are instantiated too.

\begin{lstlisting}
theorem exists_transcendental_on_circle_Qbar {r : C} (hr : r in Qbar) (hr0 : r != 0) :
 exists t : C, t != 0 and Transcendental (Qbar) t and t * conj t = r * conj r
 and t * conj t in Qbar
\end{lstlisting}
\noindent\footnotesize Uses: \texttt{conj\_mem\_Qbar}, \texttt{exists\_transcendental\_on\_circle}\normalsize


\subsection*{\texttt{mem\_Qbar\_iff} \hfill \tagpr}
\begin{lstlisting}
theorem mem_Qbar_iff {a : C} : a in Qbar <-> IsAlgebraic Q a
\end{lstlisting}
\noindent\footnotesize Uses: \texttt{Qbar}\normalsize


\end{document}
